The following definitions and lemmas are preparatory for algorithmically finding gröbner basis. The algorithm itself is discussed in \ref{subsec}

\begin{definition}
    We will write $\sPoly{f}{F}$ for the remainder on division of $f$ by the ordered s-tuple $F = (\vectorComponentsX{f}{s})$. If $F$ is a Gröbner basis for $\ideal{\vectorComponentsX{f}{s}}$, then we can regard F as a set. 
\end{definition}

Now that we have a consise way to refer to the remainder of a division by two arbitrary parameters, our dialog around the subject will be less cluttered. A deterministic approach to cancelling leading terms would help in the construction of our algorithm. Below is a definition that introduces the tool we will use to craft the algorithms kernel. 

\begin{definition}
    Let $f,g \in \polyRing$ be nonzero polynomials.
    
    \begin{enumerate}
        \item If multideg$(f) = \alpha$ and multideg$(g) = \beta$, then let $\gamma = \vectorComponentsX{\gamma}{n}$, where $\gamma_i = max(\alpha_i, \beta_i)$
        
    \end{enumerate}

\end{definition}